\section{Relation to Prior Work}

Global weak solutions trace back to Leray and Hopf \cite{Leray1934,Hopf1951}. The standard PDE framework is synthesized by Ladyzhenskaya and Temam \cite{Ladyzhenskaya1969,Temam2001}. Partial regularity and continuation criteria due to Caffarelli--Kohn--Nirenberg, Prodi, Serrin, Escauriaza--Seregin--Sver\'ak, and Koch--Tataru identify regularity mechanisms adjacent to the bridge architecture used here, while Beale--Kato--Majda remains the model derivative blow-up criterion in the inviscid setting \cite{CaffarelliKohnNirenberg1982,Prodi1959,Serrin1962,EscauriazaSereginSverak2003,KochTataru2001,BealeKatoMajda1984}. The present paper does not replace those results. It integrates them with an explicit same-surface closure route on the classical equation and discharges the scale-barrier, nonlinear-compactness, and gradient-control bridges needed for the final theorem.

The main distinctions of the present argument are explicit:
\begin{enumerate}[leftmargin=*]
\item the decisive estimates are written on the classical Euclidean equation rather than on a modified or geometric surrogate;
\item the compactness and gradient packages are tied to the same scale-barrier input and the same approximation surface;
\item the final theorem is obtained by discharging the bridge architecture rather than by appealing to a separate hidden continuation mechanism.
\end{enumerate}
